\documentclass[11pt]{article}

\usepackage[utf8]{inputenc}
\usepackage[margin=1in]{geometry}
\usepackage{amsmath,amssymb,mathtools,enumitem,fancyhdr,booktabs}
\usepackage[misc]{ifsym}
\usepackage{hyperref}

\setlength\textwidth{7in}
\setlength\parindent{0pt}
\oddsidemargin=-0.5cm
\textheight=665pt
\pagestyle{fancy}
\setlength{\headheight}{13.6pt}
\renewcommand{\headrulewidth}{0.1pt}
\renewcommand{\footrulewidth}{0.1pt}
\lhead{\scshape Pontificia Universidad Católica de Chile}
\rhead{\scshape IMC UC}
\cfoot{\thepage}

\newcommand{\norm}[1]{\left\lVert #1\right\rVert}
\newcommand{\abs}[1]{\left|#1\right|}
\def\vec   #1{\mbox{\boldmath $#1$}{}}
\def\ten   #1{\mbox{\boldmath $#1$}{}}
\def\mat   #1{\mbox{\boldmath $\mathsf #1$}}
\def\R{\mathbb R}
\def\dx{\mbox{\(\,\mathrm{d}x\)}}
\def\dX{\mbox{\(\,\mathrm{d}X\)}}
\def\dA{\mbox{\(\,\mathrm{d}A\)}}
\DeclareMathOperator{\dive}{div}
\DeclareMathOperator{\Div}{Div}
\DeclareMathOperator{\tr}{tr}
\DeclareMathOperator{\Cof}{Cof}
\DeclareMathOperator{\sym}{sym}

\AtBeginDocument{\vspace*{-3.2\baselineskip}}

\begin{document}
    \begin{flushleft}
        \small
        \vspace{1pt}
        \textbf{IMT3130} \hfill
        \textbf{23/05/2026}
    \end{flushleft}
    \begin{center}
        \LARGE\textbf{Pauta Ayudantía 9}\\
        \vspace{3pt}
        \normalsize\textbf{Mecánica del continuo: estabilidad, termodinámica y linealización}\\
        \normalsize Ayudante: Bastián Herrera \Letter{\hspace{1pt}: \texttt{biherrera@uc.cl}}
        \rule{\textwidth}{0.05mm}
    \end{center}

\section*{Problemas y soluciones}

\begin{enumerate}[label=\textbf{\arabic*.}, leftmargin=*]

    \item \textbf{Enunciado.} Sea
    \[
        V=H^1_0(\Omega)^d,
        \qquad
        Q=L^2_0(\Omega),
    \]
    y considere el problema mixto de Stokes--Brinkman: hallar $(\vec v,p)\in V\times Q$ tal que
    \[
    \begin{aligned}
        a(\vec v,\vec w)+b(\vec w,p)&=F(\vec w) &&\forall\vec w\in V,\\
        b(\vec v,q)&=0 &&\forall q\in Q,
    \end{aligned}
    \]
    donde
    \[
        a(\vec v,\vec w)=2\mu\int_\Omega D(\vec v):D(\vec w)\dx
        +\alpha\int_\Omega \vec v\cdot\vec w\dx,
        \qquad
        b(\vec w,q)=-\int_\Omega q\,\dive\vec w\dx.
    \]
    Pruebe continuidad, coercividad por Korn, unicidad del problema homogéneo, estabilidad, y discuta los casos $\alpha=0$ y $\mu=0$.

    \textbf{Solución.}
    \begin{enumerate}[label=(\alph*)]
        \item Por Cauchy--Schwarz,
        \[
        \begin{aligned}
            \abs{a(\vec v,\vec w)}
            &\le
            2\mu\norm{D(\vec v)}_{L^2}\norm{D(\vec w)}_{L^2}
            +\alpha\norm{\vec v}_{L^2}\norm{\vec w}_{L^2}\\
            &\le
            C\norm{\vec v}_{H^1}\norm{\vec w}_{H^1}.
        \end{aligned}
        \]
        En el último paso usamos que $\norm{D(\vec v)}_{L^2}\le C\norm{\vec v}_{H^1}$.
        Además,
        \[
        \begin{aligned}
            \abs{b(\vec w,q)}
            &=
            \abs{\int_\Omega q\,\dive\vec w\dx}\\
            &\le
            \norm{q}_{L^2}\norm{\dive\vec w}_{L^2}
            \le C\norm{q}_{L^2}\norm{\vec w}_{H^1}.
        \end{aligned}
        \]
        Como $\vec w\in H^1_0(\Omega)^d$, Korn da
        \[
            \norm{\vec w}_{H^1(\Omega)}\le C_K\norm{D(\vec w)}_{L^2(\Omega)}.
        \]
        Luego, usando $\mu>0$ y $\alpha\ge0$,
        \[
            a(\vec w,\vec w)
            =
            2\mu\norm{D(\vec w)}_{L^2}^2+\alpha\norm{\vec w}_{L^2}^2
            \ge c\norm{\vec w}_{H^1}^2.
        \]
        En particular, $a$ es coerciva en todo $V$, y por tanto en el núcleo solenoidal.

        \item Si $F=0$, tomamos $\vec w=\vec v$ y $q=p$. Como $b(\vec v,p)=0$,
        \[
            a(\vec v,\vec v)=0.
        \]
        Por coercividad, $\vec v=\vec0$. La primera ecuación queda
        \[
            b(\vec w,p)=0\qquad\forall \vec w\in V.
        \]
        Aplicando la inf--sup,
        \[
            \beta\norm{p}_{L^2}
            \le
            \sup_{\vec w\in V\setminus\{\vec0\}}
            \frac{\abs{b(\vec w,p)}}{\norm{\vec w}_{H^1}}
            =0.
        \]
        Así, $p=0$. La coercividad identifica la velocidad; la inf--sup identifica la presión.

        \item Para la estabilidad de $\vec v$, usamos nuevamente $\vec w=\vec v$ y $q=p$:
        \[
            a(\vec v,\vec v)=F(\vec v).
        \]
        Entonces
        \[
            c\norm{\vec v}_{H^1}^2
            \le
            \norm{F}_{V'}\norm{\vec v}_{H^1},
            \qquad
            \norm{\vec v}_{H^1}\le C\norm{F}_{V'}.
        \]
        Para $p$, de la primera ecuación,
        \[
            b(\vec w,p)=F(\vec w)-a(\vec v,\vec w).
        \]
        Por inf--sup y continuidad de $a$,
        \[
            \beta\norm{p}_{L^2}
            \le
            \norm{F}_{V'}+C\norm{\vec v}_{H^1}
            \le C\norm{F}_{V'}.
        \]
        En consecuencia,
        \[
            \norm{\vec v}_{H^1(\Omega)}+\norm{p}_{L^2(\Omega)}
            \le C\norm{F}_{V'}.
        \]

        \item Si $\alpha=0$, todavía
        \[
            a(\vec w,\vec w)=2\mu\norm{D(\vec w)}_{L^2}^2,
        \]
        y Korn entrega coercividad porque $\mu>0$. En cambio, si $\mu=0$,
        \[
            a(\vec w,\vec w)=\alpha\norm{\vec w}_{L^2}^2,
        \]
        lo que no controla $\norm{\vec w}_{H^1}$. Por eso el límite puramente Darcy no queda cubierto por esta formulación en $H^1_0(\Omega)^d$; requiere otro marco funcional, típicamente con velocidades en $H(\dive;\Omega)$.
    \end{enumerate}

    \item \textbf{Enunciado.} En $\Omega_0$, considere
    \[
        \rho_0c_v\dot\theta-\Div_{\vec X}(\ten\kappa\nabla_{\vec X}\theta)
        =
        R_0+g_{\mathrm{mec}},
        \qquad
        \ten\kappa\nabla_{\vec X}\theta\cdot\vec N=0
        \quad\text{en }\partial\Omega_0,
    \]
    con
    \[
        g_{\mathrm{mec}}
        =
        3\alpha_TK\theta_\star\,\ten F^{-\top}:\dot{\ten F}.
    \]
    Reescriba la fuente con $J=\det\ten F$, derive la formulación débil con Euler implícito, pruebe Lax--Milgram y discuta qué ocurre si la fuente contiene un término implícito $m(\vec X)\theta^{n+1}$.

    \textbf{Solución.}
    \begin{enumerate}[label=(\alph*)]
        \item Usando la fórmula de Jacobi para $J=\det\ten F$,
        \[
            \dot J=\Cof(\ten F):\dot{\ten F}
            =
            J\ten F^{-\top}:\dot{\ten F}
        \]
        implica
        \[
            \ten F^{-\top}:\dot{\ten F}=\frac{\dot J}{J}.
        \]
        Luego
        \[
            g_{\mathrm{mec}}
            =
            3\alpha_TK\theta_\star\frac{\dot J}{J}.
        \]
        Esta fuente representa una contribución proporcional a la tasa relativa de cambio de volumen.

        \item Euler implícito entrega
        \[
            \rho_0c_v\frac{\theta^{n+1}-\theta^n}{\Delta t}
            -\Div_{\vec X}(\ten\kappa\nabla_{\vec X}\theta^{n+1})
            =
            R_0^{n+1}+g_{\mathrm{mec}}^{n+1}.
        \]
        Para $q\in H^1(\Omega_0)$, multiplicamos por $q$ e integramos en $\Omega_0$:
        \[
        \begin{aligned}
            &\frac{\rho_0c_v}{\Delta t}
            \int_{\Omega_0}(\theta^{n+1}-\theta^n)q\dX
            -\int_{\Omega_0}
            \Div_{\vec X}(\ten\kappa\nabla_{\vec X}\theta^{n+1})q\dX\\
            &\qquad=
            \int_{\Omega_0}
            (R_0^{n+1}+g_{\mathrm{mec}}^{n+1})q\dX.
        \end{aligned}
        \]
        Integrando por partes el término de difusión,
        \[
        \begin{aligned}
            -\int_{\Omega_0}
            \Div_{\vec X}(\ten\kappa\nabla_{\vec X}\theta^{n+1})q\dX
            &=
            \int_{\Omega_0}
            \ten\kappa\nabla_{\vec X}\theta^{n+1}\cdot\nabla_{\vec X}q\dX\\
            &\quad
            -\int_{\partial\Omega_0}
            (\ten\kappa\nabla_{\vec X}\theta^{n+1}\cdot\vec N)q\dA\\
            &=
            \int_{\Omega_0}
            \ten\kappa\nabla_{\vec X}\theta^{n+1}\cdot\nabla_{\vec X}q\dX,
        \end{aligned}
        \]
        porque la condición de borde es adiabática. Reordenando,
        \[
        \begin{aligned}
            &\frac{\rho_0c_v}{\Delta t}\int_{\Omega_0}\theta^{n+1}q\dX
            +\int_{\Omega_0}\ten\kappa\nabla_{\vec X}\theta^{n+1}\cdot\nabla_{\vec X}q\dX\\
            &\qquad=
            \int_{\Omega_0}(R_0^{n+1}+g_{\mathrm{mec}}^{n+1})q\dX
            +\frac{\rho_0c_v}{\Delta t}\int_{\Omega_0}\theta^nq\dX.
        \end{aligned}
        \]
        El término de borde desaparece por la condición adiabática.

        \item Definimos
        \[
        \begin{aligned}
            a(\theta,q)
            &=
            \frac{\rho_0c_v}{\Delta t}\int_{\Omega_0}\theta q\dX
            +\int_{\Omega_0}\ten\kappa\nabla_{\vec X}\theta\cdot\nabla_{\vec X}q\dX,\\
            L(q)
            &=
            \int_{\Omega_0}(R_0^{n+1}+g_{\mathrm{mec}}^{n+1})q\dX
            +\frac{\rho_0c_v}{\Delta t}\int_{\Omega_0}\theta^nq\dX.
        \end{aligned}
        \]
        La forma $a$ es continua. Si $\norm{\ten\kappa}_{L^\infty}\le\kappa_1$, entonces
        \[
        \begin{aligned}
            \abs{a(\theta,q)}
            &\le
            \frac{\rho_0c_v}{\Delta t}
            \norm{\theta}_{L^2}
            \norm{q}_{L^2}
            +\kappa_1
            \norm{\nabla_{\vec X}\theta}_{L^2}
            \norm{\nabla_{\vec X}q}_{L^2}\\
            &\le
            C\norm{\theta}_{H^1(\Omega_0)}
            \norm{q}_{H^1(\Omega_0)}.
        \end{aligned}
        \]
        Además,
        \[
        \begin{aligned}
            a(\theta,\theta)
            &=
            \frac{\rho_0c_v}{\Delta t}\norm{\theta}_{L^2}^2
            +\int_{\Omega_0}\ten\kappa\nabla_{\vec X}\theta
            \cdot\nabla_{\vec X}\theta\dX\\
            &\ge
            \frac{\rho_0c_v}{\Delta t}\norm{\theta}_{L^2}^2
            +\kappa_0\norm{\nabla_{\vec X}\theta}_{L^2}^2\\
            &\ge
            \min\left\{\frac{\rho_0c_v}{\Delta t},\kappa_0\right\}
            \norm{\theta}_{H^1(\Omega_0)}^2.
        \end{aligned}
        \]
        Si
        \[
            R_0^{n+1},\quad g_{\mathrm{mec}}^{n+1},\quad \theta^n\in L^2(\Omega_0),
        \]
        entonces $L\in H^1(\Omega_0)'$, pues
        \[
        \begin{aligned}
            \abs{L(q)}
            &\le
            \left(
            \norm{R_0^{n+1}}_{L^2}
            +\norm{g_{\mathrm{mec}}^{n+1}}_{L^2}
            +\frac{\rho_0c_v}{\Delta t}\norm{\theta^n}_{L^2}
            \right)
            \norm{q}_{L^2}\\
            &\le
            C\norm{q}_{H^1(\Omega_0)}.
        \end{aligned}
        \]
        Por Lax--Milgram existe un único $\theta^{n+1}\in H^1(\Omega_0)$.

        \item Si
        \[
            g_{\mathrm{mec}}^{n+1}=m(\vec X)\theta^{n+1}+\widetilde g^{n+1},
        \]
        y se mueve el término implícito al lado izquierdo, aparece
        \[
        \begin{aligned}
            a_m(\theta,q)
            &=
            \frac{\rho_0c_v}{\Delta t}\int_{\Omega_0}\theta q\dX
            +\int_{\Omega_0}\ten\kappa\nabla_{\vec X}\theta\cdot\nabla_{\vec X}q\dX
            -\int_{\Omega_0}m\theta q\dX.
        \end{aligned}
        \]
        Una condición suficiente para conservar coercividad es
        \[
            \frac{\rho_0c_v}{\Delta t}-m(\vec X)\ge m_0>0
            \qquad\text{c.t.p. en }\Omega_0.
        \]
        Entonces
        \[
            a_m(\theta,\theta)
            \ge m_0\norm{\theta}_{L^2}^2+\kappa_0\norm{\nabla\theta}_{L^2}^2
            \ge c\norm{\theta}_{H^1}^2.
        \]
    \end{enumerate}

    \item \textbf{Enunciado.} Sea
    \[
        \vec\varphi(\vec X)=\vec X+\vec u(\vec X),
        \qquad
        \ten F=\ten I+\nabla_{\vec X}\vec u,
        \qquad
        \ten E=\frac12(\ten F^\top\ten F-\ten I),
    \]
    y la energía de Saint Venant--Kirchhoff
    \[
        \Psi(\ten E)=\frac{\lambda}{2}(\tr\ten E)^2+\mu\,\ten E:\ten E,
        \qquad
        \mu>0,\quad \lambda+\frac{2\mu}{d}>0.
    \]
    Derive $\ten S=\partial\Psi/\partial\ten E$, muestre que $\ten P=\ten F\ten S$, derive la formulación débil material, linealice alrededor de $\ten F=\ten I$ y pruebe bien planteamiento del problema linealizado por Korn y Lax--Milgram.

    \textbf{Solución.}
    \begin{enumerate}[label=(\alph*)]
        \item Derivamos primero la energía respecto de $\ten E$. Como
        \[
            \tr\ten E=E_{kk},
            \qquad
            \ten E:\ten E=E_{ij}E_{ij},
        \]
        se tiene
        \[
        \begin{aligned}
            \frac{\partial\Psi}{\partial E_{ij}}
            &=
            \lambda(\tr\ten E)\frac{\partial(\tr\ten E)}{\partial E_{ij}}
            +\mu\frac{\partial(E_{kl}E_{kl})}{\partial E_{ij}}\\
            &=
            \lambda(\tr\ten E)\delta_{ij}
            +2\mu E_{ij}.
        \end{aligned}
        \]
        Por tanto,
        \[
            \ten S=\lambda(\tr\ten E)\ten I+2\mu\ten E.
        \]
        Equivalentemente, para toda variación $\ten G$ de $\ten E$,
        \[
            D\Psi(\ten E)[\ten G]
            =
            \ten S:\ten G.
        \]
        Ahora derivamos $\Psi(\ten E(\ten F))$ respecto de $\ten F$. Si $\ten H$ es una variación de $\ten F$, entonces
        \[
            D\ten E(\ten F)[\ten H]
            =
            \frac12(\ten H^\top\ten F+\ten F^\top\ten H).
        \]
        Usando índices y que $\ten S$ es simétrica,
        \[
        \begin{aligned}
            D(\Psi\circ \ten E)(\ten F)[\ten H]
            &=
            S_{ij}\frac12(H_{ki}F_{kj}+F_{ki}H_{kj})\\
            &=
            \frac12H_{ki}F_{kj}S_{ij}
            +\frac12H_{kj}F_{ki}S_{ij}\\
            &=
            H_{ki}F_{kj}S_{ij}
            =
            (\ten F\ten S)_{ki}H_{ki}
            =
            \ten F\ten S:\ten H.
        \end{aligned}
        \]
        Luego
        \[
            \ten P=\frac{\partial(\Psi\circ\ten E)}{\partial\ten F}=\ten F\ten S.
        \]

        \item Sea
        \[
            V=\{\vec\eta\in H^1(\Omega_0)^d:\vec\eta=\vec0\text{ en }\Gamma_D\}.
        \]
        Desde
        \[
            -\Div_{\vec X}\ten P=\vec B,
            \qquad
            \ten P\vec N=\overline{\vec T}\text{ en }\Gamma_N,
        \]
        multiplicamos por $\vec\eta\in V$ e integramos:
        \[
            -\int_{\Omega_0}\Div_{\vec X}\ten P\cdot\vec\eta\dX
            =
            \int_{\Omega_0}\vec B\cdot\vec\eta\dX.
        \]
        Por integración por partes,
        \[
        \begin{aligned}
            -\int_{\Omega_0}\Div_{\vec X}\ten P\cdot\vec\eta\dX
            &=
            \int_{\Omega_0}\ten P:\nabla_{\vec X}\vec\eta\dX
            -\int_{\partial\Omega_0}(\ten P\vec N)\cdot\vec\eta\dA\\
            &=
            \int_{\Omega_0}\ten P:\nabla_{\vec X}\vec\eta\dX
            -\int_{\Gamma_N}\overline{\vec T}\cdot\vec\eta\dA.
        \end{aligned}
        \]
        En $\Gamma_D$ no aparece contribución porque $\vec\eta=\vec0$. Por tanto, se obtiene: hallar $\vec u$ admisible tal que
        \[
            \int_{\Omega_0}\ten P(\ten F(\vec u)):\nabla_{\vec X}\vec\eta\dX
            =
            \int_{\Omega_0}\vec B\cdot\vec\eta\dX
            +\int_{\Gamma_N}\overline{\vec T}\cdot\vec\eta\dA
            \qquad\forall\vec\eta\in V.
        \]
        Es no lineal porque $\ten P$ depende no linealmente de $\nabla\vec u$.

        \item Como
        \[
            \ten E
            =
            \frac12\left(\nabla\vec u+\nabla\vec u^\top+\nabla\vec u^\top\nabla\vec u\right),
        \]
        la linealización para pequeñas deformaciones elimina el término cuadrático $\nabla\vec u^\top\nabla\vec u$:
        \[
            \ten E\approx \varepsilon(\vec u)
            :=
            \frac12(\nabla\vec u+\nabla\vec u^\top).
        \]
        Entonces
        \[
            \tr\ten E
            \approx
            \tr\varepsilon(\vec u)
            =
            \dive\vec u,
            \qquad
            \ten S
            \approx
            \lambda(\dive\vec u)\ten I+2\mu\varepsilon(\vec u).
        \]
        Además, como $\ten F=\ten I+\nabla\vec u\approx\ten I$, el producto $\ten P=\ten F\ten S$ queda, a primer orden,
        \[
            \ten P
            \approx
            \ten S
            \approx
            \lambda(\dive\vec u)\ten I+2\mu\varepsilon(\vec u).
        \]
        El problema linealizado es: hallar $\vec u\in V$ tal que
        \[
            a(\vec u,\vec\eta)=L(\vec\eta)\qquad\forall\vec\eta\in V,
        \]
        con
        \[
        \begin{aligned}
            a(\vec u,\vec\eta)
            &=
            \int_{\Omega_0}
            \left[
            2\mu\,\varepsilon(\vec u):\varepsilon(\vec\eta)
            +\lambda(\dive\vec u)(\dive\vec\eta)
            \right]\dX,\\
            L(\vec\eta)
            &=
            \int_{\Omega_0}\vec B\cdot\vec\eta\dX
            +\int_{\Gamma_N}\overline{\vec T}\cdot\vec\eta\dA.
        \end{aligned}
        \]

        \item La continuidad de $a$ se obtiene de Cauchy--Schwarz:
        \[
        \begin{aligned}
            \abs{a(\vec u,\vec\eta)}
            &\le
            2\mu\norm{\varepsilon(\vec u)}_{L^2}
            \norm{\varepsilon(\vec\eta)}_{L^2}
            +|\lambda|
            \norm{\dive\vec u}_{L^2}
            \norm{\dive\vec\eta}_{L^2}\\
            &\le
            C\norm{\vec u}_{H^1(\Omega_0)}
            \norm{\vec\eta}_{H^1(\Omega_0)}.
        \end{aligned}
        \]
        Para la coercividad usamos que, para toda matriz simétrica $\ten A$,
        \[
            \ten A=\ten A_0+\frac{\tr\ten A}{d}\ten I,
            \qquad
            \tr\ten A_0=0,
            \qquad
            \abs{\ten A}^2=\abs{\ten A_0}^2+\frac{(\tr\ten A)^2}{d}.
        \]
        Entonces
        \[
        \begin{aligned}
            2\mu\,\ten A:\ten A+\lambda(\tr\ten A)^2
            &=
            2\mu\abs{\ten A_0}^2
            +\left(\lambda+\frac{2\mu}{d}\right)(\tr\ten A)^2\\
            &\ge
            c_E\left(
            \abs{\ten A_0}^2+\frac{(\tr\ten A)^2}{d}
            \right)
            =
            c_E\abs{\ten A}^2,
        \end{aligned}
        \]
        con $c_E>0$ porque $\mu>0$ y $\lambda+2\mu/d>0$. Aplicando esto a $\ten A=\varepsilon(\vec u)$,
        \[
            a(\vec u,\vec u)\ge c_E\norm{\varepsilon(\vec u)}_{L^2(\Omega_0)}^2.
        \]
        Como $|\Gamma_D|>0$, Korn implica
        \[
            \norm{\vec u}_{H^1(\Omega_0)}
            \le C_K\norm{\varepsilon(\vec u)}_{L^2(\Omega_0)}.
        \]
        Luego $a$ es coerciva en $V$. Si $\vec B\in L^2(\Omega_0)^d$ y $\overline{\vec T}\in L^2(\Gamma_N)^d$, entonces $L\in V'$ porque
        \[
        \begin{aligned}
            \abs{L(\vec\eta)}
            &\le
            \norm{\vec B}_{L^2(\Omega_0)}
            \norm{\vec\eta}_{L^2(\Omega_0)}
            +\norm{\overline{\vec T}}_{L^2(\Gamma_N)}
            \norm{\vec\eta}_{L^2(\Gamma_N)}\\
            &\le
            C\left(
            \norm{\vec B}_{L^2(\Omega_0)}
            +\norm{\overline{\vec T}}_{L^2(\Gamma_N)}
            \right)\norm{\vec\eta}_{H^1(\Omega_0)}.
        \end{aligned}
        \]
        En el último paso usamos la desigualdad de traza. Por Lax--Milgram existe una única solución del problema linealizado.
    \end{enumerate}

\end{enumerate}

\end{document}
